% xelatex graph_analyzer_flow.tex
\documentclass[11pt, a4paper]{article}

\usepackage{fontspec}
\usepackage{polyglossia}
\setdefaultlanguage{russian}
\setmainfont{Liberation Serif}
\setsansfont{Liberation Sans}
\setmonofont[Scale=0.82]{DejaVu Sans Mono}

\usepackage[
  a4paper,
  top=18mm, bottom=18mm,
  left=18mm, right=18mm
]{geometry}

\usepackage[dvipsnames, svgnames, table]{xcolor}
\usepackage{tikz}
\usetikzlibrary{
  arrows.meta,
  shapes.geometric,
  shapes.misc,
  positioning,
  fit,
  backgrounds,
  calc,
  decorations.pathreplacing,
  shadows.blur,
}

\usepackage[
  unicode, hidelinks,
  pdftitle={Поток граф-аналитика},
]{hyperref}

% ── цветовая палитра ──────────────────────────────────────────────
\definecolor{TGBlue}    {RGB}{33,  150, 243}
\definecolor{CLPurple}  {RGB}{103,  58, 183}
\definecolor{PyGold}    {RGB}{245, 158,  11}
\definecolor{LLMGreen}  {RGB}{16,  185, 129}
\definecolor{LightBg}   {RGB}{248, 249, 252}
\definecolor{LaneTG}    {RGB}{227, 242, 253}
\definecolor{LaneCL}    {RGB}{237, 231, 246}
\definecolor{LanePy}    {RGB}{255, 248, 225}
\definecolor{LaneLLM}   {RGB}{209, 250, 229}
\definecolor{ArrowGray} {RGB}{ 80,  80,  80}
\definecolor{ErrRed}    {RGB}{220,  50,  50}

% ── общие стили узлов ─────────────────────────────────────────────
\tikzset{
  every node/.style={font=\sffamily\small},
  % базовый блок
  block/.style={
    rectangle, rounded corners=4pt,
    minimum width=44mm, minimum height=9mm,
    text centered, text width=42mm,
    draw, line width=0.7pt,
    blur shadow={shadow blur steps=4, shadow xshift=0.5mm, shadow yshift=-0.5mm,
                 shadow blur radius=1.5mm, shadow opacity=25},
  },
  % Telegram
  tg/.style={block, fill=LaneTG,  draw=TGBlue!70,  text=TGBlue!90!black},
  % CL / мастер
  cl/.style={block, fill=LaneCL,  draw=CLPurple!70, text=CLPurple!90!black},
  % Python
  py/.style={block, fill=LanePy,  draw=PyGold!80,   text=PyGold!90!black},
  % LLM / внешнее
  llm/.style={block, fill=LaneLLM, draw=LLMGreen!80, text=LLMGreen!80!black},
  % ромб-решение
  decision/.style={
    diamond, aspect=2.2,
    minimum width=46mm, minimum height=12mm,
    text centered, text width=38mm,
    draw=CLPurple!70, fill=LaneCL,
    line width=0.7pt,
    font=\sffamily\footnotesize, text=CLPurple!90!black,
  },
  % маленький узел (деталь)
  sub/.style={
    rectangle, rounded corners=3pt,
    minimum width=36mm, minimum height=7.5mm,
    text centered, text width=34mm,
    draw, line width=0.5pt, dashed,
    font=\sffamily\footnotesize,
  },
  py-sub/.style={sub, fill=LanePy!60,  draw=PyGold!60,   text=PyGold!80!black},
  llm-sub/.style={sub, fill=LaneLLM!60, draw=LLMGreen!60, text=LLMGreen!70!black},
  % стрелки
  arr/.style={-{Stealth[length=5pt]}, line width=1pt, color=ArrowGray},
  arr-dashed/.style={arr, dashed, color=ArrowGray!60},
  err-arr/.style={arr, color=ErrRed, dashed},
  % метка на стрелке
  lbl/.style={font=\sffamily\scriptsize, text=ArrowGray, fill=white,
              inner sep=1.5pt, midway},
}

\pagestyle{empty}

\begin{document}

% ════════════════════════════════════════════════════════════════
\begin{center}
{\large\bfseries\sffamily Поток граф-аналитика}\\[2pt]
{\small\color{gray} от входящего файла до структурированного отчёта}
\end{center}

\vspace{4mm}

\begin{center}
\scalebox{0.72}{%
\begin{tikzpicture}[
  node distance = 7mm and 22mm,
  x=1mm, y=1mm,
]

% ── Шаг 0: пользователь ──────────────────────────────────────────
\node[tg] (user) {👤 Пользователь\\отправляет .md файл};

% ── Шаг 1: Telegram ──────────────────────────────────────────────
\node[tg, below=of user] (tg-in)
  {\textbf{Telegram}\\входящий документ};

% ── Шаг 2: мастер — fetch ────────────────────────────────────────
\node[cl, below=of tg-in] (fetch)
  {\textbf{мастер CL}\\fetch-document\\(getFile → download)};

% ── Шаг 3: решение — файл скачан? ────────────────────────────────
\node[decision, below=8mm of fetch] (ok-fetch) {Файл\\скачан?};

% ── Шаг 4: поток активируется ────────────────────────────────────
\node[cl, below=8mm of ok-fetch] (поток)
  {\textbf{поток-граф-аналитик}\\(CL)\\выполнить-с-файлом};

% ── Шаг 5: Python CLI ────────────────────────────────────────────
\node[py, below=of поток] (cli)
  {\textbf{cli.py}\\python3 graph\_analyzer/cli.py\\graph.md [models.py]};

% ── Шаг 6: парсер (два подузла рядом) ────────────────────────────
\node[py-sub, below=8mm of cli, xshift=-24mm] (parser-g)
  {parser.py\\parse\_mermaid()\\→ nodes, edges};
\node[py-sub, below=8mm of cli, xshift=+24mm] (parser-m)
  {parser.py\\parse\_models\_simple()\\→ [{name, fields}]};

% соединитель
\coordinate (split) at ($(cli.south)+(0,-4mm)$);

% ── Шаг 7: аналитик ──────────────────────────────────────────────
\node[py, below=26mm of cli] (analyzer)
  {\textbf{analyzer.py}\\analyze() — 3 шага};

% ── Три LLM-запроса ──────────────────────────────────────────────
\node[llm-sub, below=8mm of analyzer, xshift=-34mm] (llm1)
  {① Компенсация\\compensation\_gap};
\node[llm-sub, below=8mm of analyzer] (llm2)
  {② Крайние случаи\\edge\_cases};
\node[llm-sub, below=8mm of analyzer, xshift=+34mm] (llm3)
  {③ Достижимость\\reachability};

% LLM API
\node[llm, below=26mm of analyzer] (openrouter)
  {\textbf{OpenRouter API}\\anthropic/claude-opus-4-5\\→ findings JSON};

% ── Форматтер ────────────────────────────────────────────────────
\node[py, below=of openrouter] (fmt)
  {\textbf{formatter.py}\\format\_report()\\→ текст отчёта};

% ── мастер возвращает ─────────────────────────────────────────────
\node[cl, below=of fmt] (reply)
  {\textbf{мастер CL}\\send-message\\усечь до 3800 симв.};

% ── Telegram → пользователь ──────────────────────────────────────
\node[tg, below=of reply] (tg-out)
  {\textbf{Telegram}\\исходящее сообщение};
\node[tg, below=of tg-out] (user-out)
  {👤 Пользователь\\получает отчёт};

% ── error path (справа от ok-fetch) ──────────────────────────────
\node[tg, right=22mm of ok-fetch, yshift=4mm,
      fill=red!10, draw=ErrRed!70, text=ErrRed,
      minimum width=34mm, text width=32mm] (err)
  {send-message\\«Не удалось скачать»};

% ════ СТРЕЛКИ ═════════════════════════════════════════════════════

% user → tg-in → fetch → ok-fetch
\draw[arr] (user)     -- (tg-in);
\draw[arr] (tg-in)    -- (fetch);
\draw[arr] (fetch)    -- (ok-fetch);

% ok-fetch → Да → поток
\draw[arr] (ok-fetch) -- node[lbl, right=1pt]{Да} (поток);

% ok-fetch → Нет → err
\draw[err-arr] (ok-fetch) -- node[lbl, above=1pt]{Нет} (err);

% поток → cli
\draw[arr] (поток)    -- (cli);

% cli → два парсера
\draw[arr] (split) -- (parser-g);
\draw[arr] (split) -- (parser-m);
\draw[arr-dashed] (cli.south) -- (split);

% парсеры → analyzer (сливаются)
\coordinate (merge) at ($(analyzer.north)+(0,4mm)$);
\draw[arr] (parser-g) |- (merge);
\draw[arr] (parser-m) |- (merge);
\draw[arr-dashed] (merge) -- (analyzer.north);

% analyzer → три LLM-запроса
\coordinate (asplit) at ($(analyzer.south)+(0,-4mm)$);
\draw[arr-dashed] (analyzer.south) -- (asplit);
\draw[arr] (asplit) -- (llm1);
\draw[arr] (asplit) -- (llm2);
\draw[arr] (asplit) -- (llm3);

% три запроса → openrouter (сливаются)
\coordinate (amerge) at ($(openrouter.north)+(0,4mm)$);
\draw[arr] (llm1) |- (amerge);
\draw[arr] (llm2) |- (amerge);
\draw[arr] (llm3) |- (amerge);
\draw[arr-dashed] (amerge) -- (openrouter.north);

% openrouter → fmt → reply → tg-out → user-out
\draw[arr] (openrouter) -- node[lbl, right=1pt]{findings JSON} (fmt);
\draw[arr] (fmt)         -- (reply);
\draw[arr] (reply)       -- (tg-out);
\draw[arr] (tg-out)      -- (user-out);

% ════ ФОНОВЫЕ ПОЛОСЫ (lanes) ══════════════════════════════════════
\begin{scope}[on background layer]

  % Telegram lane (левый бордюр)
  \node[fit=(user)(tg-in)(tg-out)(user-out),
        fill=TGBlue!5, draw=TGBlue!25, rounded corners=6pt,
        inner sep=5pt, label={[font=\sffamily\tiny\color{TGBlue}]above left:Telegram}] {};

  % CL lane
  \node[fit=(fetch)(ok-fetch)(поток)(reply),
        fill=CLPurple!5, draw=CLPurple!20, rounded corners=6pt,
        inner sep=5pt, label={[font=\sffamily\tiny\color{CLPurple}]above left:CL / мастер}] {};

  % Python lane
  \node[fit=(cli)(parser-g)(parser-m)(analyzer)(fmt),
        fill=PyGold!5, draw=PyGold!25, rounded corners=6pt,
        inner sep=5pt, label={[font=\sffamily\tiny\color{PyGold}]above left:Python / graph\_analyzer}] {};

  % LLM lane
  \node[fit=(llm1)(llm2)(llm3)(openrouter),
        fill=LLMGreen!5, draw=LLMGreen!25, rounded corners=6pt,
        inner sep=5pt, label={[font=\sffamily\tiny\color{LLMGreen!70!black}]above right:OpenRouter LLM}] {};

\end{scope}

% ════ ЛЕГЕНДА ═════════════════════════════════════════════════════
\node[below=14mm of user-out, xshift=32mm,
      draw=gray!40, fill=LightBg, rounded corners=4pt,
      inner sep=5pt, font=\sffamily\scriptsize, text=gray!80] (legend) {
  \begin{tabular}{cl}
    \textcolor{TGBlue}{\rule{8pt}{5pt}} & Telegram \\[2pt]
    \textcolor{CLPurple}{\rule{8pt}{5pt}} & Common Lisp / мастер \\[2pt]
    \textcolor{PyGold}{\rule{8pt}{5pt}} & Python / graph\_analyzer \\[2pt]
    \textcolor{LLMGreen}{\rule{8pt}{5pt}} & OpenRouter API (LLM) \\[2pt]
    \textcolor{ErrRed}{- - -} & Путь ошибки \\
  \end{tabular}
};

\end{tikzpicture}%
}% scalebox
\end{center}

\vfill
\begin{center}
{\footnotesize\color{gray}
  \texttt{мастер/потоки/граф-аналитик.lisp} · \texttt{tools/graph\_analyzer/} ·
  development\_system
}
\end{center}

\end{document}
